\chapter{Open Problems and Research Directions}
\label{ch:open-problems}

\abstract*{Post-quantum cryptography is a field in rapid motion. This final chapter identifies concrete open problems and research directions, with particular attention to opportunities where a physics background provides a decisive advantage. From quantum cryptanalysis of lattice problems to hardware optimization and fundamental complexity questions, the field needs researchers who understand both quantum mechanics and cryptographic security.}

% =============================================================
% SOURCE: notes.tex §6.1, §6.3 (lines 7877–9409)
% REUSE: ~50%. Add concrete numbered open problems,
%        recent developments 2024-2025.
% =============================================================

\section{Quantum Cryptanalysis}
\label{sec:quantum-cryptanalysis}

% TODO: Adapt from notes.tex §6.1.1
% - Quantum algorithms for lattice problems: current state
% - Quantum speedups for lattice sieving and enumeration
% - Quantum walks on lattice graphs
% - The holy grail: is there a polynomial-time quantum algorithm for SVP?

\section{Physical Side-Channel Research}
\label{sec:physical-research}

% TODO: Adapt from notes.tex §6.1.2
% - EM analysis of PQC implementations
% - Deep learning for side-channel attacks
% - Countermeasure design: physics-informed approaches
% - Quantum side channels: a new frontier

\section{Hardware Optimization}
\label{sec:hardware}

% TODO: Adapt from notes.tex §6.1.3
% - FPGA and ASIC implementations of NTT
% - Gaussian sampling hardware
% - Energy-efficient PQC for IoT
% - Photonic computing for lattice operations

\section{Cryptographic Agility}
\label{sec:crypto-agility}

% TODO: Adapt from notes.tex §6.3.1
% - Design principles for algorithm-agnostic systems
% - Protocol negotiation and version management
% - The cost of agility vs the cost of replacement
% - Lessons from the SHA-1 and MD5 transitions

\section{Preparing for New Breaks}
\label{sec:preparing-breaks}

% TODO: Adapt from notes.tex §6.3.2
% - What if Module-LWE is broken?
% - Defense in depth: diversification of assumptions
% - Monitoring the cryptanalytic landscape
% - The role of crypto competitions and challenges

\section{The 50-Year Horizon}
\label{sec:horizon}

% TODO: Adapt from notes.tex §6.3.3
% - Quantum internet and its cryptographic implications
% - Post-quantum + quantum: the full security stack
% - Information-theoretic vs computational security in the long term
% - New computational models: quantum advantage beyond Shor

\section{Concrete Open Problems}
\label{sec:open-problem-list}

% TODO: New section — numbered list for reference by PhD students
% Example format:
% \begin{enumerate}
%   \item \textbf{Quantum SVP:} Is there a polynomial-time quantum algorithm ...
%   \item \textbf{Tight reductions:} Can the security loss in ... be eliminated?
%   \item ...
% \end{enumerate}

%% Exercises
\begin{prob}
\label{prob:ch15-agility}
Design a cryptographic agility framework for a web service currently using TLS~1.3 with X25519 and Ed25519. Your framework must allow swapping to any NIST-standardized PQC algorithm within 48 hours of a break announcement. Identify the components that need abstraction and the operational procedures required.
\end{prob}

\begin{prob}
\label{prob:ch15-research}
Choose one of the open problems listed in Section~\ref{sec:open-problem-list}. Write a 2-page research proposal outlining: (a)~the problem's significance, (b)~your proposed approach, (c)~what tools or techniques from physics you would bring to bear, and (d)~what a successful outcome would look like.
\end{prob}
